\section{The Prayer}

The Regina Coeli is both an Easter Marian antiphon and, in an expanded form, a pious exercise of the Latin Church prayed in place of the Angelus during Easter Time. Its repeated alleluia joins Mary's joy to the Resurrection of the Son she bore. The form below includes the four-line antiphon, its versicle and response, and the customary collect.

This bilingual text reproduces Dom Gaspar Lefebvre's 1925 \textit{Daily Missal with Vespers for Sundays and Feasts}, printed pp.~13--14. The spelling, capitalization, punctuation, stress accents, and historical English have not been modernized or translated by this project. The source begins the prayer at the foot of p.~13 and continues it on p.~14.

\bilingualprayer{The Regina Coeli}{Received Latin and historical English from Lefebvre's 1925 \textit{Daily Missal}, pp.~13--14. The source's historical season, posture, and indulgence rubric is discussed separately and is not a statement of current discipline.}{%
Regína coeli laetáre, allelúia;\par
Quia quem meruísti portáre, allelúia:\par
Resurréxit sicut dixit, allelúia:\par
Ora pro nobis Deum, allelúia.\par
\smallskip
\VR{Gaude et laetáre, Virgo María, allelúia.}{Quia surréxit Dóminus vere, allelúia.}
\textbf{Orémus.}\par
Deus, qui per resurrectiónem Fílii tui Dómini nostri Jesu Christi mundum laetificáre dignátus es: praesta, quaesumus; ut per ejus Genitrícem Vírginem Maríam, perpétuae capiámus gáudia vitae. Per eúmdem Christum Dóminum nostrum.\par
\textbf{R.} Amen.}{%
Joy to thee, O Queen of heaven, alleluia!\par
He Whom thou wast meet to bear, alleluia,\par
As He promis'd, hath arisen, alleluia;\par
Pour for us to Him thy prayer, alleluia.\par
\smallskip
\VR{Rejoice and be glad, O Virgin Mary, alleluia.}{For the Lord hath risen indeed, alleluia.}
\textbf{Let us pray.}\par
O God, Who didst vouchsafe to give joy to the world through the resurrection of Thy Son our Lord Jesus Christ; grant, we beseech Thee, that through His Mother, the Virgin Mary, we may obtain the joys of everlasting life. Through the same Christ our Lord.\par
\textbf{R.} Amen.}

One may pray the Latin, the historical English, or another received form devotionally; the prayer need not be doubled. A vernacular form used for an indulgenced recitation requires approval by competent ecclesiastical authority, and a liturgical celebration follows its governing book. This witness spells the title \textit{Regina coeli}; the 1999 \textit{Enchiridion} prints \textit{Regina caeli}, modern \textit{Iesu} and \textit{eius}, and a shorter \textit{Per Christum} conclusion. Those current forms are not silently substituted for this 1925 witness.
